%% ma: L11-B06-0028
%% lop: 11
%% chuong: 2
%% bai: 6
%% dang: 11.6.12
%% loai: TL
%% muc: TH
%% dap_an: "-"
%% nguon: "de-on-tap-giua-ki-1-toan-11-nam-2026-2027-truong-thpt-ngo-quyen-da-nang.pdf (M09, Đề ôn tập giữa kì 1 Toán 11 THPT Ngô Quyền, Đà Nẵng), Mã đề 0114, Phần IV Câu 2"
%% kien_thuc: "Bài 6 (cấp số cộng)"
%% trang_thai: chinh-thuc
%% ghi_chu: "Đã duyệt 29/10/2026."
\begin{ex}
Một công ty trả lương cho anh A theo phương thức sau: Mức lương quý đầu tiên là $4{,}5$ triệu đồng/quý (một năm chia thành $4$ quý). Kể từ quý tiếp theo, mỗi quý được tăng thêm $0{,}3$ triệu đồng. Hỏi tổng số tiền lương anh A nhận được sau ba năm làm việc là bao nhiêu?
\loigiai{Lương các quý lập thành cấp số cộng $u_1=4{,}5$, $d=0{,}3$. Ba năm có $12$ quý: $S_{12}=\dfrac{12}2(2\cdot4{,}5+11\cdot0{,}3)=6\cdot12{,}3=73{,}8$ (triệu đồng).}
\end{ex}
